
\subsubsection{Spectral Analysis of Neural Networks}

Martin and Mahoney (2019) showed that heavy-tailed eigenvalue distributions emerge during training and correlate with generalization. Their alpha exponent captures power-law decay, though alpha failed to discriminate model quality in prior attempts on this research program (p = 0.55). The WeightWatcher tool \citep{Martin2021WeightWatcher} operationalized this framework for computing spectral metrics across layers. These methods compute features directly but do not examine the variance of randomized estimators.

Li et al. (2018) estimated intrinsic dimensionality of objective landscapes using random subspace training. Their participation ratio (PR = ($\Sigma \lambda )^2 / \Sigma \lambda ^{2})$ captures effective rank. The present work adopts PR but focuses on its coefficient of variation across SVD seeds.

Randomized SVD algorithms \citep{Halko2011Finding} underpin the extraction methodology. By projecting onto random subspaces, these algorithms trade exactness for speed. The projection introduces variance that is typically treated as noise to minimize; this study treats it as a potential signal.

\subsubsection{Model Quality Prediction}

Unterthiner et al. (2020) achieved $R^2$ > 0.98 for predicting model accuracy from weights using simple statistics aggregated across layers. Their approach is architecture-agnostic. The present work tests whether spectral variance adds interpretable signal beyond direct spectral features.

\subsubsection{Architecture-Aware Analysis}

The timm library \citep{Wightman2019timm} provides access to pretrained models with documented accuracy. Model soups \citep{Wortsman2022ModelSoups} and Git Re-Basin \citep{Ainsworth2023GitReBasin} explore weight-space relationships but focus on interpolation and alignment rather than spectral variance. Architecture-family effects (BatchNorm, GroupNorm, LayerNorm creating different spectral signatures) remain underexplored.
